\scanpage{28}
\mspage{9}
Posons
\[
\text{(2.5)}\qquad\underline S=\operatorname{gr}_{\mathcal J}\underline O_{X'}
=\coprod_{k\in\underline{\underline N}}\mathcal J^k/\mathcal J^{k+1}
\]
\[
\text{(2.6)}\qquad\operatorname{gr}_{\mathcal J}(\underline F)
=\coprod_{k\in\underline{\underline N}}\mathcal J^k\underline F/\mathcal J^{k+1}\underline F
\]
\[
\text{(2.7)}\qquad\underline K^i=R^if_*(\operatorname{gr}_{\mathcal J}(\underline F))
=\coprod_{k\in\underline{\underline N}}R^if_*(\mathcal J^k\underline F/\mathcal J^{k+1}\underline F),
\]
il est clair que $\underline K^i$ est muni d'une structure de $\underline S$-Module
gradué.

\uline{Théorème 2.1.} Supposons que $\underline K^i$ soit un $\underline S$-Module gradué
\uline{de type fini} pour $i=n-1,i=n,i=n+1$, alors:

(i) $R^nf_*(F)$ est cohérent

(ii) l'homomorphisme $\unsure{\psi|\varphi}_n$, (2.3), est un isomorphisme
topo\inserted{l}ogique\inserted{, La filtration naturelle du deuxième membre
de (2.3.) est $\mathcal J$-bonne.}

(iii) le système projectif des $R^nf_*(F_k)$ vérifie
la condition de Mittag-Leffler uniforme.

La démonstration est très facile à partir de
E.G.A.$0_{\mathrm{III}}$.13.7.7.\inserted{(cf. E.G.A. III 3.4.2.)}, à condition de rectifier comme suit
le texte page 78:

\uline{Supprimer}

\begin{tabular}{ll}
ligne 6 & `` et $(R^{n+1}T(A_k))_{k\in\underline{\underline Z}}$ ''\\[1ex]
ligne 13 & `` et $p+q=n+1$''\\[1ex]
ligne 23 & `` et pour $n+1$''.
\end{tabular}

\uline{Théorème 2.2.} Soit $A$ un anneau adique noethérien et soit
$I$ un idéal de définition de $A$.\inserted{Soit $T$ une partie fermée de $X'=\operatorname{Spec}(A)$.}
Supposons que $I$ soit engendré par un $t\in A$. Reprenons les notations (1.31)\inserted{,}(1.32) et
(1.33). Soit $\underline F$ un $\underline O_{\widehat X}$-Module cohérent. Posons
\[
\text{\unsure{(2.8)|(2.89)}}\qquad F_0=\underline F/\mathcal J\underline F,
\]
où $\mathcal J=t\underline O_{\widehat X}$ est un idéal de définition de $\widehat X$.
\margin{27}

